% Lab 3 --- compact version (full docs: docs/web/docs/labs/lab3/)
\section{Lab 3: Incident Response and Log Forensics --- HELIX Systems}

Lab~3 inverts the student's role: instead of attacking, they act as a security
analyst responding to an active incident at \textit{HELIX Systems}. The lab
environment comes pre-staged with evidence of a simulated intrusion --- brute-force
SSH attempts, a successful login, privilege escalation via \texttt{sudo}
misconfiguration, backdoor account creation, and lateral movement to a database
server. Students must reconstruct the attack timeline from system logs
(\texttt{auth.log}, \texttt{bash\_history}) and identify the attacker's actions
and impact, following a methodology aligned with NIST SP~800-61~\cite{nist_sp80061}.
A quick teacher reference --- account table, attack timeline, and startup checklist --- is compiled in Appendix~D;
the full evidence walkthrough, forensic guide, and deployment instructions are at \path{docs/web/docs/labs/lab3/}.


\begin{figure}[H]
\centering
\includegraphics[width=0.95\textwidth]{lab3-topology.png}
\caption{Lab~3 network topology}
\label{fig:lab3-topology}
\end{figure}

\subsubsection*{Network Addressing}

\begin{table}[H]
\centering
\resizebox{\textwidth}{!}{\begin{tabular}{@{}lllll@{}}
\toprule
\textbf{Node} & \textbf{Role} & \textbf{Interface} & \textbf{IP} & \textbf{Network} \\
\midrule
Parrot-OS    & Analyst station   & eth0  & 192.168.0.x (DHCP) & 192.168.0.0/24 (Homelab)       \\
pfSense-LAB3 & Firewall/DNAT     & WAN   & 192.168.0.x (DHCP) & 192.168.0.0/24                 \\
pfSense-LAB3 & Firewall          & LAN   & 172.16.2.1         & 172.16.2.0/30                  \\
VyOS-LAB3    & Internal router   & eth0  & 172.16.2.2         & 172.16.2.0/30                  \\
VyOS-LAB3    & Internal router   & eth1  & 192.168.50.1       & 192.168.50.0/24 (Net-Servers)  \\
VyOS-LAB3    & Internal router   & eth2  & 192.168.60.1       & 192.168.60.0/24 (Net-Internal) \\
Server-Web   & Compromised web   & eth0  & 192.168.50.10      & 192.168.50.0/24                \\
Server-DB    & Internal DB       & eth0  & 192.168.60.10      & 192.168.60.0/24                \\
\bottomrule
\end{tabular}}
\caption{Lab~3 IP addressing scheme}
\label{tab:lab3-ip}
\end{table}

\subsubsection*{MITRE ATT\&CK Coverage}

\begin{table}[H]
\centering
\small
\begin{tabularx}{\textwidth}{llX}
\toprule
\textbf{MITRE ID} & \textbf{Technique} & \textbf{Evidence in lab} \\
\midrule
T1110.001 & Brute Force: Password Guessing & \texttt{auth.log} failed attempts \\
T1078     & Valid Accounts                 & \texttt{devops} account compromise \\
T1548.003 & Abuse Elevation Control Mechanism: Sudo and Sudo Caching          & NOPASSWD rule for \texttt{find} in \texttt{sudoers.d/devops} \\
T1136.001 & Create Account: Local          & Backdoor \texttt{maint} account \\
T1021.004 & Remote Services: SSH           & Lateral movement to Server-DB \\
T1005     & Data from Local System         & \texttt{/opt/helix/data/} exfil marker \\
\bottomrule
\end{tabularx}
\caption{Lab~3 MITRE ATT\&CK technique coverage}
\label{tab:lab3-mitre}
\end{table}

\newpage
\subsubsection*{Forensic Findings}

\begin{table}[H]
\centering
\small
\resizebox{\textwidth}{!}{\begin{tabular}{@{}lll@{}}
\toprule
\textbf{Finding} & \textbf{Artefact / Location} & \textbf{Host} \\
\midrule
Brute-force source (\texttt{185.220.101.47}) & \texttt{/var/log/auth.log} --- 25 failed SSH for \texttt{devops} & Server-Web \\
Initial compromise (03:17)   & \texttt{auth.log} successful login as \texttt{devops}                                        & Server-Web \\
Privilege escalation (03:18) & \texttt{sudo /usr/bin/find} NOPASSWD; \texttt{auth.log} + \texttt{/etc/sudoers.d/devops}    & Server-Web \\
Backdoor persistence (03:19) & \texttt{useradd maint}; \texttt{/etc/passwd} + \texttt{/etc/sudoers.d/maint}                 & Server-Web \\
Attacker commands            & \texttt{/home/devops/.bash\_history}                                                         & Server-Web \\
Lateral movement (03:21)     & \texttt{auth.log} \texttt{devops} SSH from \texttt{192.168.50.10}                            & Server-DB  \\
Data exfiltration (03:22)    & \texttt{/opt/helix/data/.exfil\_marker} + \texttt{clients.db}                               & Server-DB  \\
\bottomrule
\end{tabular}}
\caption{Lab~3 forensic findings and evidence locations (all timestamps UTC as recorded in \texttt{/var/log/auth.log})}
\label{tab:lab3-findings}
\end{table}

\subsection{Problems Encountered and Solutions}

\subsubsection{LVM Volume Group Naming Conflict}

When pre-staging Server-Web's evidence files from the EVE-NG host, the same
volume group conflict from Lab~2
recurred: both the host and the guest disk used \texttt{ubuntu-vg}. The resolution
procedure was identical --- renaming the guest VG to \texttt{helix-web-vg} before
activation --- and the naming rule was formally added to the project's infrastructure
runbook to prevent recurrence in future labs.

\subsubsection{Read-only NBD Device}

On one occasion, the QCOW2 image mounted via \texttt{qemu-nbd} was presented as
read-only despite the absence of the \texttt{--read-only} flag. Investigation revealed that a
previous mount session had not been cleanly disconnected (\texttt{qemu-nbd~-{}-disconnect}), leaving the kernel's NBD client in an inconsistent state.
The fix was to explicitly disconnect any stale NBD device before re-mounting, and
to add a pre-mount check to the project's image manipulation scripts.

\subsubsection{Analyst Group Membership for Log Access}

The student account (\texttt{analyst}) was created with restricted permissions so
that it could not read system logs directly without elevated group membership. However, \texttt{/var/log/auth.log}
requires membership of the \texttt{adm} group on Ubuntu. The pre-staging script was
updated to add \texttt{analyst} to \texttt{adm} at image preparation time, preserving
the educational constraint (no \texttt{sudo} to \texttt{root}) while allowing access
to the evidence files that are the focus of the lab.

\subsubsection{pfSense Console Access via QEMU Monitor}

The pfSense serial console was inaccessible through the standard EVE-NG Telnet
interface due to a configuration mismatch between the QEMU serial port definition
and pfSense's expected console device (\texttt{ttyu0} vs \texttt{ttyS0}). Console
access was recovered by connecting directly to the QEMU monitor socket and issuing
a \texttt{sendkey} sequence to navigate the pfSense boot menu, where the console
port assignment could be corrected persistently.
